\documentclass[10pt,aspectratio=169]{beamer}
% Preámbulo modular para presentaciones Beamer (Modelación Estadística)
% Diseño y paleta institucional del Tecnológico de Monterrey

\documentclass[aspectratio=169,xcolor={dvipsnames,table}]{beamer}

\usepackage[utf8]{inputenc}
\usepackage[spanish,mexico]{babel}
\usepackage{amsmath,amssymb,amsfonts,mathtools}
\usepackage{graphicx}
\usepackage{listings}
\usepackage{booktabs}
\usepackage{tikz}
\usetikzlibrary{matrix,arrows,decorations.pathmorphing,shapes.geometric,positioning}

% Paleta Institucional Tecnológico de Monterrey
\definecolor{TecAzul}{HTML}{0039A6}       % Azul TEC: color primario institucional
\definecolor{TecAzulOscuro}{HTML}{1A2E51} % Azul marino oscuro
\definecolor{TecRojo}{HTML}{EC2661}       % Rojo de acento (paleta miTec)
\definecolor{TecGris}{HTML}{646464}       % Gris neutro para comentarios y subtítulos
\definecolor{TecFondoBloque}{HTML}{F4F6F9}% Fondo suave para bloques

% Configuración de Tema Beamer
\usetheme{Boadilla}
\usecolortheme[named=TecAzulOscuro]{structure}
\setbeamercolor{alerted text}{fg=TecRojo}
\setbeamercolor{palette primary}{bg=TecAzulOscuro,fg=white}
\setbeamercolor{palette secondary}{bg=TecAzul,fg=white}
\setbeamercolor{palette tertiary}{bg=TecAzulOscuro!80!black,fg=white}
\setbeamercolor{title}{fg=TecAzulOscuro,bg=TecFondoBloque}
\setbeamercolor{frametitle}{fg=TecAzulOscuro,bg=TecFondoBloque}

% Bloques personalizados elegantes
\setbeamercolor{block title}{bg=TecAzulOscuro,fg=white}
\setbeamercolor{block body}{bg=TecFondoBloque,fg=black}
\setbeamercolor{block title alerted}{bg=TecRojo,fg=white}
\setbeamercolor{block body alerted}{bg=TecRojo!8,fg=black}
\setbeamercolor{block title example}{bg=TecAzul,fg=white}
\setbeamercolor{block body example}{bg=TecAzul!8,fg=black}

% Redondear bordes y añadir sombra sutil
\setbeamertemplate{blocks}[rounded][shadow=true]
\setbeamertemplate{navigation symbols}{}

% Pie de página limpio con número de diapositiva
\setbeamertemplate{footline}{
  \leavevmode%
  \hbox{%
  \begin{beamercolorbox}[wd=.7\paperwidth,ht=2.25ex,dp=1ex,left]{author in head/foot}%
    \hspace*{2ex}\insertshortauthor~~(\insertshortinstitute) \hfill \insertshorttitle
  \end{beamercolorbox}%
  \begin{beamercolorbox}[wd=.3\paperwidth,ht=2.25ex,dp=1ex,right]{date in head/foot}%
    \insertshortdate{}\hspace*{2em}
    \insertframenumber{} / \inserttotalframenumber\hspace*{2ex}
  \end{beamercolorbox}}%
  \vskip0pt%
}

% Configuración de Listings de Python para visualización pedagógica de código
\lstset{
    language=Python,
    basicstyle=\ttfamily\footnotesize,
    keywordstyle=\color{TecAzulOscuro}\bfseries,
    stringstyle=\color{TecRojo},
    commentstyle=\color{TecGris}\itshape,
    showstringspaces=false,
    breaklines=true,
    frame=lines,
    rulecolor=\color{TecAzulOscuro!30},
    backgroundcolor=\color{TecFondoBloque},
    aboveskip=0.5em,
    belowskip=0.5em,
    literate={á}{{\'a}}1 {é}{{\'e}}1 {í}{{\'i}}1 {ó}{{\'o}}1 {ú}{{\'u}}1
             {Á}{{\'A}}1 {É}{{\'E}}1 {Í}{{\'I}}1 {Ó}{{\'O}}1 {Ú}{{\'U}}1
             {ñ}{{\~n}}1 {Ñ}{{\~N}}1 {¿}{{?`}}1 {¡}{{!`}}1
}

% Metadatos institucionales por defecto
\title[Modelación Estadística]{Modelación Estadística para Ciencia de Datos e Ingeniería}
\author[J. Castillo Colmenares]{Juliho David Castillo Colmenares}
\institute[Tec de Monterrey]{Tecnológico de Monterrey \\ Escuela de Ingeniería y Ciencias}
\date{\today}
% Comandos y macros estadísticas compartidas para las presentaciones Beamer (Metropolis)

% Conjuntos numéricos básicos
\newcommand{\R}{\mathbb{R}}
\newcommand{\N}{\mathbb{N}}
\newcommand{\Q}{\mathbb{Q}}
\newcommand{\Z}{\mathbb{Z}}
\newcommand{\set}[1]{\left\{ #1 \right\}}
\newcommand{\abs}[1]{\left|#1\right|}
\newcommand{\norm}[1]{\left\| #1 \right\|}

% Comandos estadísticos y probabilísticos del curso
\newcommand{\Var}{\operatorname{Var}}
\newcommand{\cov}{\operatorname{Cov}}
\newcommand{\corr}{\rho}
\newcommand{\comb}[2]{\begin{pmatrix} #1 \\ #2 \end{pmatrix}}
\newcommand{\s}{\sigma}
\newcommand{\particion}{\mathfrak{P}}
\newcommand{\card}[1]{\#\left( #1 \right)}
\newcommand{\seq}[3]{\set{#1}_{#2}^{#3}}
\newcommand{\E}{\mathbb{E}}
\newcommand{\Prob}{\mathbb{P}}

% Entornos pedagógicos y cajas en español académico natural
\newenvironment{discusion}{%
  \begin{alertblock}{Pregunta de Discusión}%
}{%
  \end{alertblock}%
}

\newenvironment{reflexion}{%
  \begin{alertblock}{Reflexión para el Aula}%
}{%
  \end{alertblock}%
}

\newenvironment{paradoja}{%
  \begin{alertblock}{Paradoja de Exploración}%
}{%
  \end{alertblock}%
}

\newenvironment{motivacion}{%
  \begin{exampleblock}{Motivación: Ciencia de Datos en la Práctica}%
}{%
  \end{exampleblock}%
}

\newenvironment{retoclase}{%
  \begin{block}{Reto Rápido en Equipo}%
}{%
  \end{block}%
}

\title[Guía de pruebas de hipótesis]{Sección 07.05: Guía paso a paso para una prueba de hipótesis}\subtitle{De la pregunta al reporte}
\author[J. Castillo Colmenares]{Juliho Castillo Colmenares}\institute{Tecnológico de Monterrey}\date{}
\begin{document}
\begin{frame}[plain]\titlepage\end{frame}
\begin{frame}{Hoja de ruta}\small\begin{enumerate}\item Formular hipótesis.\item Muestrear y estandarizar.\item Comparar, decidir y comunicar.\end{enumerate}\end{frame}
\begin{frame}{Fuente y alcance}\small Esta guía resume los pasos de la nota y los organiza como una secuencia reproducible.\begin{block}{Pregunta guía}¿Qué decisiones deben quedar documentadas?\end{block}\end{frame}
\begin{frame}{Paso 1: hipótesis nula}\small La hipótesis nula expresa la afirmación de referencia sobre el parámetro y contiene el valor $A_0$.\[H_0:A=A_0.\]\end{frame}
\begin{frame}{Paso 1: alternativa}\small La alternativa expresa la desviación de interés: $A<A_0$, $A>A_0$ o $A\neq A_0$. Su dirección define las colas de la prueba.\end{frame}
\begin{frame}{Paso 2: muestra}\small Tome una muestra aleatoria de unidades u ocurrencias y calcule el estimador correspondiente, por ejemplo una media $A_m$.\begin{alertblock}{Calidad} La independencia y el plan de muestreo sostienen la interpretación.\end{alertblock}\end{frame}
\begin{frame}{Paso 2: estadístico}\small Resuma la diferencia entre $A_m$ y $A_0$ mediante un error estándar. Para una media con $\sigma$ conocida, $Z=(A_m-A_0)/(\sigma/\sqrt n)$.\end{frame}
\begin{frame}{Paso 3: valor normal}\small La nota introduce el valor normal estándar $Z$ y su probabilidad asociada. La distribución bajo $H_0$ determina el área relevante.\end{frame}
\begin{frame}{Paso 3: niveles y colas}\small El nivel de significación se fija antes de observar los datos. Una alternativa unilateral usa una cola; una bilateral reparte $\alpha$ en dos.\end{frame}
\begin{frame}{Paso 4: comparación}\small Compare la probabilidad asociada al estadístico con $\alpha$.\[p\leq\alpha\Rightarrow\text{rechazo};\qquad p>\alpha\Rightarrow\text{no rechazo}.\]\end{frame}
\begin{frame}{Paso 5: decisión}\small La decisión es condicional al modelo y al nivel escogido. No rechazar no significa aceptar de forma definitiva la hipótesis nula.\end{frame}
\begin{frame}{Paso 6: comunicación}\small Reporte la hipótesis, la estadística, el nivel, el valor $p$ y una conclusión en el contexto de la variable medida.\end{frame}
\begin{frame}{Aplicación: restaurante}\small Un restaurante afirma un tiempo medio de entrega de $20$ minutos con desviación estándar de $3$ minutos. Se desea evaluar la afirmación con una muestra.\end{frame}
\begin{frame}{Aplicación: planteamiento}\small\[H_0:\mu=20.\]La alternativa debe expresar la preocupación concreta: $\mu>20$, $\mu<20$ o $\mu\neq20$.\end{frame}
\begin{frame}{Aplicación: estandarización}\small Con media muestral $\bar x$ y tamaño $n$,\[Z=\frac{\bar x-20}{3/\sqrt n}.\]El signo indica la dirección de la evidencia.\end{frame}
\begin{frame}{Aplicación: área}\small La probabilidad asociada con $Z$ se obtiene de la normal estándar y se convierte en valor $p$ según la alternativa elegida.\end{frame}
\begin{frame}{Aplicación: decisión}\small Compare el valor $p$ con $\alpha$. La conclusión debe indicar si la evidencia es suficiente para cuestionar la afirmación de veinte minutos.\end{frame}
\begin{frame}{Errores frecuentes}\small\begin{itemize}\item Intercambiar $H_0$ y $H_a$.\item Elegir la cola después de calcular.\item Confundir no rechazo con aceptación.\item Omitir unidades y contexto.\end{itemize}\end{frame}
\begin{frame}{Lista de comprobación}\small\begin{enumerate}\item Parámetro identificado.\item Hipótesis completas.\item Supuestos revisados.\item $\alpha$ fijado.\item Estadística y valor $p$ reportados.\end{enumerate}\end{frame}
\begin{frame}{Plantilla de reporte}\small\begin{block}{Redacción} ``Para $H_0:\cdots$ frente a $H_a:\cdots$, se obtuvo $T=\cdots$ y $p=\cdots$. Al nivel $\alpha=\cdots$, [rechazamos/no rechazamos] $H_0$.''\end{block}\end{frame}
\begin{frame}{Síntesis}\small La secuencia completa transforma una pregunta aplicada en una decisión estadística auditable y comunicable.\end{frame}
\begin{frame}{Conexión con el capítulo}\small Los pasos se reutilizan en pruebas $t$, proporciones, varianzas, homogeneidad y varias proporciones.\end{frame}
\end{document}
